% ----------------------------------------------------------------------
%              Latex PhD template for the University of Deusto
% ----------------------------------------------------------------------


%: Style file for Latex
% Most style definitions are in the external file PhDthesisPSnPDF.
% In this template package, it can be found in ./Latex/Classes/
\documentclass[twoside,12pt]{Latex/Classes/PhDthesisPSnPDF}

%\usepackage[lined, boxed]{algorithm2e} % for algorithm figures creation.

%: Macro file for Latex
% Macros help you summarise frequently repeated Latex commands.
% Here, they are placed in an external file /Latex/Macros/MacroFile1.tex
% An macro that you may use frequently is the figuremacro (see introduction.tex)
%---------------------------------------------------------------
% Macros
% version 3 by Igor Ruiz-Agundez 2011
% version 2 by Jakob Suckale 2007
% version 1 by Harish Bhanderi 2002
%---------------------------------------------------------------

% This file contains macros that can be called up from connected TeX files
% It helps to summarise repeated code, e.g. figure insertion (see below).



%---------------------------------------------------------------
% Figures
%---------------------------------------------------------------


% Makes the \InsertFig macro compatible both with one or two columns
\makeatletter
\newlength \figwidth
\if@twocolumn
  \setlength \figwidth {\columnwidth}
\else
  \setlength \figwidth {\textwidth}
\fi
\makeatother

% \InsertFig allows inserting figures
% Parameters
% 1 --> Filename
% 2 --> Label for referencing
% 3 --> Title describing the figure (caption)
% 4 --> Description of the figure
% 5 --> Figure width, range [0,1]. If parameter is left blank the figure size is not change
% 6 --> Any other option for \includegraphics
% Usage:
% \InsertFig{}{}{}{}{}{}
%
\newcommand{\InsertFig}[6]{%
	\ifthenelse{\isempty{#5}}%
	{% if #1 is empty
		\begin{figure}[htbp!]
		\centering
		\includegraphics[#6]{#1}%
		\caption{#3}{\textbf{#4}}
		\label{#2}
		\end{figure}    
	}
	{% if #1 is not empty
		\begin{figure}[htbp!]
		\centering
		\includegraphics[width=#5\figwidth,#6]{#1}%
		\caption{#3}{\textbf{#4}}
		\label{#2}
		\end{figure}
	}
}

%% Simple version of \InsertFig
%\newcommand{\InsertFig}[5]{
%  \begin{figure}[htbp]
%   	\centering
%    \includegraphics[width=#4\textwidth,#5]{#1}%
%    \caption{#3}
%    \label{#2}
%  \end{figure}
%}



% insert a centered figure with caption
% parameters 1:filename, 2:label, 3:title, 
\newcommand{\figuremacro}[3]{
	\begin{figure}[htbp]
		\centering
		\includegraphics[width=1\textwidth]{#1}
		\caption[#3]{\textbf{#3}}
		\label{#2}
	\end{figure}
}


% insert a centered figure with caption and description
% parameters 1:filename, 2:label, 3:title, 4:description
\newcommand{\figuremacroD}[4]{
	\begin{figure}[htbp]
		\centering
		\includegraphics[width=1\textwidth]{#1}
		\caption[#3]{\textbf{#3} - #4}
		\label{#2}
	\end{figure}
}

% insert a centered figure with caption and description AND WIDTH
% parameters 1:filename, 2:label, 3:title, 4:description, 5: textwidth
% textwidth 1 means as text, 0.5 means half the width of the text
\newcommand{\figuremacroDW}[5]{
	\begin{figure}[htbp]
		\centering
		\includegraphics[width=#5\textwidth]{#1}
		\caption[#3]{\textbf{#3} - #4}
		\label{#2}
	\end{figure}
}

% inserts a figure with wrapped around text; only suitable for NARROW figs
% o is for outside on a double paged document; others: l, r, i(inside)
% text and figure will each be half of the document width
% note: long captions often crash with adjacent content; take care
% in general: above 2 macro produce more reliable layout
\newcommand{\figuremacroN}[3]{
	\begin{wrapfigure}{o}{0.5\textwidth}
		\centering
		\includegraphics[width=0.48\textwidth]{#1}
		\caption[#2]{{\small\textbf{#2} - #3}}
		\label{#1}
	\end{wrapfigure}
}




% Estas definiciones son para el comando \InsertFigBox
\newlength{\anchoFigura}
\newlength{\anchoFloat}
\addtolength{\fboxsep}{2\fboxsep}
%\renewcommand{\capfont}{\normalfont\normalcolor\sffamily\small}
%\renewcommand{\caplabelfont}{\normalfont\normalcolor\sffamily\bfseries\small}

% El comando \InsertFigBox nos permite insertar figuras en un marco
% Los parametros son:
% 1 --> Fichero de la imagen
% 2 --> Etiqueta (label) para referencias
% 3 --> Texto a pie de imagen
% 4 -> Porcentaje del ancho de página que ocupará la figura (de 0 a 1)
% 5 --> Opciones que queramos pasarle al \includegraphics
\newcommand{\InsertFigBox}[5]{%
  \setlength{\anchoFloat}{#4\textwidth}%
  \addtolength{\anchoFloat}{-4\fboxsep}%
  \setlength{\anchoFigura}{\anchoFloat}%
  \begin{figure}%
    \begin{center}%
      \Ovalbox{%
        \begin{minipage}{\anchoFloat}%
          \begin{center}%
            \includegraphics[width=\anchoFigura,#5]{#1}%
            \caption{#3}%
            \label{#2}%
          \end{center}%
        \end{minipage}
      }%
    \end{center}%
  \end{figure}%
}



%---------------------------------------------------------------
% Misc
%---------------------------------------------------------------

% predefined commands by Harish
\newcommand{\PdfPsText}[2]{
  \ifpdf
     #1
  \else
     #2
  \fi
}


%---------------------------------------------------------------
% Locales
%---------------------------------------------------------------


%%
%% Para quitar traducciones raras (Cuadros)
%% Ha de usarse cada vez que se seleccione el idioma
%%
\newcommand{\MejorarTraducciones}{%
	   \renewcommand{\contentsname}{Tabla de contenido}
	   \renewcommand{\listfigurename}{Lista de figuras}
	   \renewcommand{\listtablename}{Lista de tablas}
	   \renewcommand{\tablename}{Tabla}
	   \renewcommand{\lstlistingname}{Lista}
	   \renewcommand{\nomname}{Acrónimos}
	   \renewcommand{\sectionname}{Sección}
	   \renewcommand{\acknowledgementstitle}{Agradecimientos}
}%

%%
%% Permite referenciar figuras directamente de la forma \figurename~\ref{label}
%%
%%
\newcommand{\fig}[1]{\mbox{\MakeLowercase\figurename~\ref{#1}}}

%%
%% Permite referenciar secciones directamente de la forma Sección \ref{label}
%%
%%
\newcommand{\sect}[1]{\mbox{\MakeLowercase\sectionname~\ref{#1}}}

%%
%% Permite referenciar secciones directamente de la forma Capítulo \ref{label}
%%
%%
\newcommand{\chap}[1]{\mbox{\MakeLowercase\chaptername~\ref{#1}}}

\newcommand{\Chap}[1]{\mbox{\chaptername~\ref{#1}}}

%%
%% Permite referenciar secciones directamente de la forma Capítulo \ref{label}
%%
%%
\newcommand{\tabl}[1]{\mbox{\MakeLowercase\tablename~\ref{#1}}}

%:--------------------- TODO command --------------------------------
\newcommand{\todo}[1]{\textcolor{red}{@TODO: #1}}


%---------------------------------------------------------------
% Source code
%---------------------------------------------------------------


%%
%% Para escribir extractos de codigo
%%
%% Las tabulaciones se substituyen por dos espacios
%\fvset{tabsize=2}
%% Creamos un nuevo environment de fancyvrb para los ejemplos enmarcados
%\DefineVerbatimEnvironment{VerbEj}{BVerbatim}{fontsize=\small,samepage=true,commandchars=\\\{\}}
%% Colo de fondo
%\definecolor{grisfondo}{gray}{0.9}
%% Environment para extractos de codigo
%\newenvironment{codigo}%
%{\VerbatimEnvironment\begin{Sbox}\begin{VerbEj}}%
%{\end{VerbEj}\end{Sbox}\setlength{\fboxsep}{8pt}\begin{center}\fcolorbox{black}{grisfondo}{\TheSbox}\end{center}}
%
%% Otro formato más bonito para código fuente
%\newcommand{\codigofuente}[3]{%
%  \lstinputlisting[language=#1,caption={#2}]{#3}%
%}






%: ----------------------------------------------------------------------
%:                  TITLE PAGE: name, degree,..
% ----------------------------------------------------------------------
% below is to generate the title page with crest and author name

% Title of the dissertation
\newcommand{\thesistitle}{Eine neue Bestimmung der Moleküldimensionen:
A New Determination of Molecular Dimensions}

%Author name
\newcommand{\authorname}{Albert Einstein}

%First advisor name
\newcommand{\advisorname}{Dr. Alfred Kleiner}

% if output to PDF then put the following in PDF header
\ifpdf
    \hypersetup{pdfauthor={\authorname},%
            pdftitle={\thesistitle},%
            pdfsubject={PhD Thesis},%
            pdfkeywords={keyword1, keyword2, keyword3},%
            pdfproducer={LaTeX},%
            pdfcreator={pdfLaTeX}
}
\fi

% ----------------------------------------------------------------------
% This section below defines front covert (external and internal)
% Shield logo
\crest{\includegraphics[width=2cm]{Deusto_Shield}}
% Full logo
%\crest{\includegraphics[width=6cm]{UDeusto}}
\university{Universidad de Deusto}
\degree{Tesis doctoral presentada por }
\author{\authorname} 
\collegeordept{dentro del Programa de Doctorado en Física Experimental}
\textadvisor{Dirigida por el }
\advisor{\advisorname}
\advisortwo{}
\textsignaturecandidate{El doctorando}
\textsignatureadvisor{El director}
\cityofbirth{Bilbao}
\degreedate{enero de 2013}

% ----------------------------------------------------------------------
       
% turn of those nasty overfull and underfull hboxes
\hbadness=10000
\hfuzz=50pt

%trying to avoid orphan and widow lines
\widowpenalty10000
\clubpenalty10000

% hyphenation
%\hyphenation{so-me-thing}

%: --------------------------------------------------------------
%:                  FRONT MATTER: dedications, abstract,..
% --------------------------------------------------------------

\begin{document}

\selectlanguage{spanish}
\MejorarTraducciones

% sets line spacing
%\onehalfspacing

% Watermark
%\watermark{DRAFT	DRAFT	DRAFT	DRAFT	DRAFT	DRAFT	DRAFT	DRAFT	DRAFT}


%: ----------------------- generate cover page ------------------------

\maketitle  % command to print the title page with above variables

% Title back
% Thesis Titleback ---------------------------------------------------

\thispagestyle{empty}

\hfill

\vfill

\medskip


\noindent
\textit{\thesistitle}
\\

\noindent
Autor: \authorname

\noindent
Director: \advisorname


\vfill

\vfill

\noindent
%The following web-page address contains up to date information about this dissertation and related topics: \\
%\url{http://paginaspersonales.deusto.es/Name/}
%Las figuras de este libro han sido realizadas con imágenes de dominio público obtenidas de la Open Clip Art Library (http://openclipart.org/).

\noindent
Edited with \LaTeX2e using \textit{TeX Gyre Heros} for titles and
\textit{Charter BT} for the general content.
\\

\noindent
Printed in Bilbao

% TODO final date
\noindent
First edition, 
% Moth and year
\monthname \ \the\year

\vspace{1cm}
\hrule
\bigskip

% \cleardoublepage command ends the current page and causes all figures and tables that have so far appeared in the input to be printed. In a two-sided printing style, it also makes the next page a right-hand (odd-numbered) page, producing a blank page if necessary. 
\cleardoublepage

%%: ----------------------- cover page back side ------------------------
%% Your research institution may require reviewer names, etc.
%% This cover back side is required by Dresden Med Fac; uncomment if needed.
%
%\newpage
%\vspace{10mm}
%1. Reviewer: Name
%
%\vspace{10mm}
%2. Reviewer: 
%
%\vspace{20mm}
%Day of the defense:
%
%\vspace{20mm}
%\hspace{70mm}Signature from head of PhD committee:
%
%
%\cleardoublepage

% ----------------------------------------------------------------------



%: ----------------------- abstract ------------------------

% Your institution may have specific regulations if you need an abstract and
% where it is to be placed in the document. The default here is just after
% title.

% The original template provides and abstractseparate environment, if your
% institution requires them to be separate. I think it's easier to print the
% abstract from the complete thesis by restricting printing to the relevant
% page.
% \begin{abstractseparate}
%   
\begin{resumen} 
\selectlanguage{spanish} 
Lorem ipsum dolor sit amet, consectetur adipiscing elit. Morbi consectetur enim
vitae urna fermentum non congue leo molestie. Maecenas adipiscing consequat
augue. Duis ante ante, porttitor quis cursus vitae, cursus quis mauris. Maecenas
sed erat et massa pellentesque malesuada ullamcorper sed enim. Cras in massa non
magna fermentum pulvinar non vitae nisi. Proin mi quam, posuere ut sagittis in,
bibendum ac mauris. Fusce hendrerit nulla id tellus laoreet lobortis. Nam et
erat vitae mi pharetra bibendum quis ac enim. Duis et urna eget purus viverra
varius. Quisque rhoncus nunc sit amet diam fringilla id fringilla sem bibendum.
Duis et diam sed nulla consectetur varius. Morbi auctor nulla non turpis
condimentum nec pellentesque metus molestie. Aenean diam tellus, posuere sed
aliquam vel, pharetra ut mi. Vestibulum pharetra turpis a nisl convallis eu
gravida sapien rhoncus. Aliquam suscipit elementum lacus ornare convallis. Cras
gravida diam sed urna aliquet nec ultrices turpis aliquam.

Mauris rutrum, tortor ut auctor rutrum, enim sapien varis mauris, vel ornare
erat lectus pharetra lectus. Donec sed nisl id enim aliquam bibendum eleifend in
massa. Praesent luctus enim ac dolor eleifend iaculis. Cras imperdiet, ligula
quis consequat interdum, sem lectus pretium diam, eu dictum neque mauris vel
ligula. Integer volutpat erat sit amet tortor dictum consequat. Integer felis
urna, aliquet et dignissim eget, consectetur ut leo. In hac habitasse platea
dictumst. Sed ornare lorem sed nisl sodales adipiscing. Pellentesque habitant
morbi tristique senectus et netus et malesuada fames ac turpis egestas. Mauris
at ipsum nisi. Morbi nulla nibh, tempor eget vulputate quis, viverra ac tortor.
Proin adipiscing, dui ac pellentesque dictum, lacus tellus placerat ante, id
vulputate ipsum dolor id nisi. Integer sed felis eu neque facilisis sodales eget
vel justo. Fusce non lacus libero. Morbi porttitor, tellus quis volutpat
iaculis, lectus lorem posuere risus, a posuere massa nunc eu dui. Quisque
feugiat augue ac tellus viverra sed accumsan nisi auctor.
\end{resumen}

% \end{abstractseparate}


%: ----------------------- tie in front matter ------------------------

% The frontmatter text starts here
\frontmatter

% Thesis Dedictation ---------------------------------------------------

\begin{dedication} %this creates the heading for the dedication page

\textit{Para alguien.}

\end{dedication}

% ----------------------------------------------------------------------

\begin{resumen} 
\selectlanguage{spanish} 
Lorem ipsum dolor sit amet, consectetur adipiscing elit. Morbi consectetur enim
vitae urna fermentum non congue leo molestie. Maecenas adipiscing consequat
augue. Duis ante ante, porttitor quis cursus vitae, cursus quis mauris. Maecenas
sed erat et massa pellentesque malesuada ullamcorper sed enim. Cras in massa non
magna fermentum pulvinar non vitae nisi. Proin mi quam, posuere ut sagittis in,
bibendum ac mauris. Fusce hendrerit nulla id tellus laoreet lobortis. Nam et
erat vitae mi pharetra bibendum quis ac enim. Duis et urna eget purus viverra
varius. Quisque rhoncus nunc sit amet diam fringilla id fringilla sem bibendum.
Duis et diam sed nulla consectetur varius. Morbi auctor nulla non turpis
condimentum nec pellentesque metus molestie. Aenean diam tellus, posuere sed
aliquam vel, pharetra ut mi. Vestibulum pharetra turpis a nisl convallis eu
gravida sapien rhoncus. Aliquam suscipit elementum lacus ornare convallis. Cras
gravida diam sed urna aliquet nec ultrices turpis aliquam.

Mauris rutrum, tortor ut auctor rutrum, enim sapien varis mauris, vel ornare
erat lectus pharetra lectus. Donec sed nisl id enim aliquam bibendum eleifend in
massa. Praesent luctus enim ac dolor eleifend iaculis. Cras imperdiet, ligula
quis consequat interdum, sem lectus pretium diam, eu dictum neque mauris vel
ligula. Integer volutpat erat sit amet tortor dictum consequat. Integer felis
urna, aliquet et dignissim eget, consectetur ut leo. In hac habitasse platea
dictumst. Sed ornare lorem sed nisl sodales adipiscing. Pellentesque habitant
morbi tristique senectus et netus et malesuada fames ac turpis egestas. Mauris
at ipsum nisi. Morbi nulla nibh, tempor eget vulputate quis, viverra ac tortor.
Proin adipiscing, dui ac pellentesque dictum, lacus tellus placerat ante, id
vulputate ipsum dolor id nisi. Integer sed felis eu neque facilisis sodales eget
vel justo. Fusce non lacus libero. Morbi porttitor, tellus quis volutpat
iaculis, lectus lorem posuere risus, a posuere massa nunc eu dui. Quisque
feugiat augue ac tellus viverra sed accumsan nisi auctor.
\end{resumen}

% Thesis Acknowledgements ------------------------------------------------

% Opening of the acknowledgements

%Sort version
%this creates the heading for the acknowlegments
\begin{acknowledgements}      
%Long version
%uncommenting this line, gives a different acknowledgements heading
%\begin{acknowledgementslong} 

Aliquam pretium consequat fringilla. Duis varius sollicitudin massa vel
accumsan. Fusce id molestie ante. Pellentesque ultrices nisi sed mi semper
volutpat. Vivamus facilisis ipsum eget metus consequat porttitor. Nam ut ligula
leo, ac vestibulum nisl. Cras interdum sodales nibh sed mollis. In id mauris
dui, sit amet viverra felis. Aliquam erat volutpat. Proin tempus nisi
sollicitudin lacus posuere auctor. Etiam ut velit et tellus feugiat tempus vel
et elit. Ut vel turpis felis, vitae semper felis. Pellentesque laoreet fringilla
lacus, vel eleifend neque consectetur ac. Nam fermentum bibendum bibendum. Morbi
in nulla non odio tincidunt mattis et ac velit.

Vivamus non nisi ut nunc venenatis laoreet. Aliquam dolor nulla, sollicitudin
quis porttitor commodo, dictum at ligula. Nulla est massa, laoreet in placerat
ac, suscipit quis metus. Aenean ipsum ante, tincidunt vitae sodales a, pulvinar
et quam. Aenean erat lacus, congue commodo pharetra placerat, faucibus tincidunt
arcu. Cras mattis lectus a ante placerat eget fringilla est consectetur. Donec
consequat luctus mi. Phasellus imperdiet, arcu ut ornare venenatis, leo neque
ultrices leo, at imperdiet ante purus vitae erat.

\begin{flushright}
\textit{Thank you,}

Your name

% Moth and year
Someplace, \monthname{} \the\year

% Signature figure

%\begin{figure}[htbp!]
%\end{figure}
%\includegraphics{signature}%

\end{flushright}

%Closing of the acknowledgements 
%Sort version
\end{acknowledgements}
% Long version
%\end{acknowledgementslong}

% ------------------------------------------------------------------------




% As abstract contains various languages we set the main language again
\selectlanguage{spanish}
\MejorarTraducciones

%: ----------------------- contents ------------------------

\setcounter{secnumdepth}{5} % organisational level that receives a numbers
\setcounter{tocdepth}{5}    % print table of contents for level 3

%You can also add extra lines to the ToC or to force extra unnumbered section
%headings to be included. For example, if you wanted to add an entry called
%Preface, and you didn't want the Preface to be numbered, you'd use these
%commands:
%\ subsection*{Preface}
%\addcontentsline{toc}{subsection}{Preface} 

\cleardoublepage
\pdfbookmark[0]{\contentsname}{toc}
\tableofcontents            % print the table of contents
							% levels are:
							% 0 - chapter
							% 1 - section
							% 2 - subsection
							% 3 - subsection
%: ----------------------- list of figures/tables ------------------------

\listoffigures	% print list of figures
\listoftables  % print list of tables

%: ----------------------- glossary ------------------------

% Tie in external source file for definitions: /0_frontmatter/glossary.tex
% Glossary entries can also be defined in the main text. See glossary.tex
\include{0_frontmatter/glossary}

%\begin{multicols}{2} % \begin{multicols}{#columns}[header text][space]
%\begin{footnotesize} % scriptsize(7) < footnotesize(8) < small (9) < normal (10)

\printnomenclature[1.5cm] % [] = distance between entry and description
\label{sec:glossary} % target name for links to glossary

%\end{footnotesize}
%\end{multicols}

%defined commands
\newcommand{\edistancia}[0]{\emph{distancia estimada}}

%: --------------------------------------------------------------
%:                  MAIN DOCUMENT SECTION
% --------------------------------------------------------------

% the main text starts here with the introduction, 1st chapter,...
\mainmatter

%\renewcommand{\chaptername}{} % uncomment to print only "1" not "Chapter 1"

%: ----------------------- subdocuments ------------------------

% Parts of the thesis are included below. Rename the files as required. But take
% care that the paths match. You can also change the order of appearance by
% moving the include commands.


% this file is called up by thesis.tex
% content in this file will be fed into the main document

%: ----------------------- Introduction file header -----------------------


\begin{savequote}[50mm]
Un viaje de mil millas comienza
con el primer paso.
\qauthor{\textsc{Lao-tsé} (570 a.C.- 490 a.C.)}
\end{savequote}

\chapter{Introducción}
% the code below specifies where the figures are stored
\ifpdf
    \graphicspath{{1_introduction/figures/PNG/}{1_introduction/figures/PDF/}{1_introduction/figures/}}
\else
    \graphicspath{{1_introduction/figures/EPS/}{1_introduction/figures/}}
\fi


%------------------------------------------------------------------------- 

Como puede deducirse del trabajo de \cite{Bailey} y, posteriormente del de
\cite{Ta}... 

% ----------------------------------------------------------------------


\section{Motivación}
\label{sec:motivaton}
Curabitur placerat auctor mattis. Aliquam at augue sem. Ut adipiscing aliquet
eleifend. Nullam vitae eros vitae felis consequat placerat. Ut viverra blandit
velit at facilisis. Ut mattis pretium vestibulum. Maecenas elementum, nibh
porttitor sodales pulvinar, diam arcu molestie nunc, nec auctor magna tortor
eget eros. Nam consectetur nisl at dolor volutpat feugiat. 

Como se observa en la \fig{fig:atom} \ldots

\InsertFig{atom}{fig:atom}{Esto es un átomo}{}{0.25}{}

Mauris fringilla convallis facilisis. Quisque condimentum, enim in gravida
convallis, justo felis suscipit nisi, ac sagittis turpis nibh in sem. Curabitur
lacinia ante non lacus iaculis in lacinia turpis consectetur. Nunc sollicitudin
metus sit amet libero accumsan at rhoncus turpis volutpat. Aenean nec velit at
arcu porttitor accumsan nec et neque. In quis odio nunc. Pellentesque habitant
morbi tristique senectus et netus et malesuada fames ac turpis egestas.	

% this file is called up by thesis.tex
% content in this file will be fed into the main document

%: ----------------------- Introduction file header -----------------------


\chapter{Otro Capítulo}
\label{cha:otro_capitulo}

% the code below specifies where the figures are stored
\ifpdf
    \graphicspath{{2_another_chapter/figures/PNG/}{2_another_chapter/figures/PDF/}{2_another_chapter/figures/}}
\else
    \graphicspath{{2_another_chapter/figures/EPS/}{2_another_chapter/figures/}}
\fi


%------------------------------------------------------------------------- 

Lorem ipsum dolor sit amet, consectetur adipiscing elit. Aenean placerat, tortor
quis adipiscing malesuada, velit tellus viverra ligula, non fermentum nunc erat
fermentum metus. Nulla et dapibus erat. Proin accumsan tristique odio, vel
fermentum leo porttitor in. Aliquam eu urna turpis. Suspendisse ac eleifend
enim. Vivamus a magna at felis consequat auctor vitae vel erat. Donec faucibus
justo vel odio pulvinar id mattis arcu vulputate. Cras risus nisi, imperdiet nec
hendrerit quis, mattis eu velit. Nunc sit amet ligula metus. Morbi magna enim,
adipiscing mollis interdum quis, bibendum id lectus. 

Como se muestra en la \tabl{tab:example} \ldots

\begin{table}
\center
\caption{Una tabla de ejemplo}
\begin{tabular}{|r|l|}
  \hline
  7C0 & hexadecimal \\
  3700 & octal \\ \cline{2-2}
  11111000000 & binary \\
  \hline \hline
  1984 & decimal \\
  \hline
\end{tabular}
\label{tab:example}
\end{table}

Sed adipiscing justo fermentum tortor laoreet sed vestibulum tortor mattis.
Vivamus a enim augue. Donec nunc metus, facilisis vitae accumsan et,
pellentesque non lorem. Curabitur porta iaculis diam vitae varius. Duis at diam
in urna fringilla egestas a ac felis. Aliquam semper malesuada orci, id vehicula
mauris varius.


% ----------------------------------------------------------------------
	

% --------------------------------------------------------------
%:                  BACK MATTER: appendices, refs,..
% --------------------------------------------------------------

% the back matter: appendix and references close the thesis
\backmatter

%: ----------------------- appendix ------------------------

\appendix

\chapter{Apéndice A}
\label{chap:appendix_publications}

% the code below specifies where the figures are stored
\ifpdf
    \graphicspath{{appendix_publications/figures/PNG/}{appendix_publications/figures/PDF/}{appendix_publications/figures/}}
\else
    \graphicspath{{appendix_publications/figures/EPS/}{appendix_publications/figures/}}
\fi


%-------------------------------------------------------------------------

Suspendisse potenti. In hac habitasse platea dictumst. Nullam mattis metus eu
quam dictum blandit. Nam nisl nibh, mattis eget laoreet at, tincidunt quis
turpis. Nulla ac augue in lectus suscipit ultrices id ac odio. Ut a purus eget
nulla adipiscing dapibus a vitae mi. Vestibulum tempor odio a ipsum sollicitudin
porttitor. Phasellus porta metus a sem pellentesque eget vulputate urna
interdum. Nam euismod odio at libero tempus id eleifend ipsum viverra. In hac
habitasse platea dictumst.

% ------------------------------------------------------------------------


%: ----------------------- bibliography ------------------------

% The section below defines how references are listed and formatted The default
% below is 2 columns, small font, complete author names. Entries are also linked
% back to the page number in the text and to external URL if provided in the
% BibTex file.


% Original version:

% PhDbiblio-url2 = names small caps, title bold & hyperlinked, link to page 
%\begin{multicols}{2} % \begin{multicols}{ # columns}[ header text][ space]
%\begin{tiny} % tiny(5) < scriptsize(7) < footnotesize(8) < small (9)
%
%\bibliographystyle{Latex/Classes/PhDbiblio-url2} % Title is link if provided
%\renewcommand{\bibname}{References} % changes the header; default: Bibliography
%
%\bibliography{9_backmatter/references} % adjust this to fit your BibTex file
%
%\end{tiny}
%\end{multicols}

% Show all bibliography entries
%\nocite*


% If we want bibliography backreference, use unsrt first and the desidered one after

%\bibliographystyle{unsrt} % Defines the bibliography style

%\bibliographystyle{alpha} % Defines the bibliography style
\bibliographystyle{Latex/StyleBST/TeXiS}

%\bibliographystyle{apa-good} % Defines the bibliography style
%\bibliographystyle{natbib} % Defines the bibliography style

%\bibliographystyle{plainurl}

%\renewcommand{\bibname}{References} % changes the header; default: Bibliography

%To include the references/works cited/bibliography in your Table of Contents,
%right before the bibliography command, use the command
%\addcontentsline{toc}{section}{References}

\bibliography{9_backmatter/references} % adjust this to fit your BibTex file

% --------------------------------------------------------------
% Various bibliography styles exit. Replace above style as desired.

% in-text refs: (1) (1; 2)
% ref list: alphabetical; author(s) in small caps; initials last name; page(s)
%\bibliographystyle{Latex/Classes/PhDbiblio-case} % title forced lower case
%\bibliographystyle{Latex/Classes/PhDbiblio-bold} % title as in bibtex but bold
%\bibliographystyle{Latex/Classes/PhDbiblio-url} % bold + www link if provided

%\bibliographystyle{Latex/Classes/jmb} % calls style file jmb.bst
% in-text refs: author (year) without brackets
% ref list: alphabetical; author(s) in normal font; last name, initials; page(s)

%\bibliographystyle{plainnat} % calls style file plainnat.bst
% in-text refs: author (year) without brackets
% (this works with package natbib)

% --------------------------------------------------------------


%: back pages

% this file is called up by thesis.tex
% content in this file will be fed into the main document

%: ----------------------- introduction file header -----------------------


% the code below specifies where the figures are stored
\ifpdf
    \graphicspath{{9_backmatter/figures/PNG/}{9_backmatter/figures/PDF/}{9_backmatter/figures/}}
\else
    \graphicspath{{9_backmatter/figures/EPS/}{9_backmatter/figures/}}
\fi

% Thesis statement of originality -------------------------------------

% Depending on the regulations of your faculty you may need a declaration like
% the one below. This specific one is from the medical faculty of the university
% of Dresden.
% 
% \begin{declaration}  %this creates the heading for the declaration page
%

% I herewith declare that I have produced this work without the prohibited
% assistance of third parties and without making use of aids other than those
% specified; notions taken over directly or indirectly from other sources have
% been identified as such. This work has not previously been presented in
% identical or similar form to any examination board.

% The dissertation work was conducted from 20XX to \the\year \ under the
% supervision of Name Surname and Name Surname at the University of Deusto.

% \vspace{10mm}
% 
% City
% Bilbao,
% 
% % Signature figure
% 
% %\begin{flushright}
% %\begin{figure}[htbp!]
% %\end{figure}
% %\includegraphics{signature}%
% %\end{flushright}
% 
% 
% \end{declaration}



%-------------------------------------------------------------------------

% Blank page
\clearpage
\thispagestyle{empty}
\mbox{}

%-------------------------------------------------------------------------

\clearpage
\thispagestyle{empty}

\hfill
\vfill
\medskip
\begin{center}
\noindent
Esta tesis fue terminada de escribir en Bilbao el \today 
\end{center}

\vfill
\medskip
\vspace{1cm}
\bigskip

%-------------------------------------------------------------------------

% Blank page
\clearpage
\thispagestyle{empty}
\mbox{}

%-------------------------------------------------------------------------


\end{document}
